%%%%%%%% Manuscript v2 (v1: git commit 8fb0368) %%%%%%%%%%%%%%%%%%%%%%%%%%%%%%%%%%%%%
% TARGET_YEAR      = 2027  (ICML 2027; no official 2027 style published as of 2026-09-26)
% TEMPLATE_YEAR    = 2026  (official ICML 2026 style kit, review/anonymous mode, unmodified)
% SUBMISSION_READY = false
% Build: scripts/build_paper.sh (copies the official kit, sha256-checked, into paper/build/;
%        the style files are not vendored in the repository).
% All numbers in the text come from paper/generated/numbers.tex; all tables and the figure
% from paper/generated/*.tex, produced by scripts/make_paper_assets.py from results/raw.
%%%%%%%%%%%%%%%%%%%%%%%%%%%%%%%%%%%%%%%%%%%%%%%%%%%%%%%%%%%%%%%%%%%%%%%%%%%%%%%%%%%%%%

\documentclass{article}

\usepackage{microtype}
\usepackage{graphicx}
\usepackage{subcaption}
\usepackage{booktabs}
\usepackage{hyperref}

\newcommand{\theHalgorithm}{\arabic{algorithm}}

% Initial blind version submitted for review (official ICML 2026 style; see header):
\usepackage{icml2026}

\usepackage{amsmath}
\usepackage{amssymb}
\usepackage{mathtools}
\usepackage{amsthm}
\usepackage[capitalize,noabbrev]{cleveref}
\usepackage{pgfplots}
\usepgfplotslibrary{groupplots}
\pgfplotsset{compat=1.17}

\theoremstyle{plain}
\newtheorem{theorem}{Theorem}[section]
\newtheorem{proposition}[theorem]{Proposition}
\newtheorem{lemma}[theorem]{Lemma}
\theoremstyle{definition}
\newtheorem{definition}[theorem]{Definition}
\newtheorem{assumption}[theorem]{Assumption}
\theoremstyle{remark}
\newtheorem{remark}[theorem]{Remark}

% Visible marker for evidence that does not exist yet (never a placeholder number).
\newcommand{\todo}[1]{\textcolor{red}{\textbf{[TODO: #1]}}}
\newcommand{\notrun}{\textsc{not run}}

% AUTO-GENERATED by scripts/make_paper_assets.py from results/raw -- do not edit by hand.
\newcommand{\AdapterParams}{\ensuremath{23{,}397}}
\newcommand{\AdapterStep}{\ensuremath{0.031}}
\newcommand{\CompErrEight}{\ensuremath{0.026}}
\newcommand{\CompErrOne}{\ensuremath{0.19}}
\newcommand{\CompInv}{\ensuremath{1.7\times10^{-15}}}
\newcommand{\CompRK}{\ensuremath{6.1\times10^{-15}}}
\newcommand{\CompSeconds}{\ensuremath{0.11}}
\newcommand{\CompSlope}{\ensuremath{-0.95}}
\newcommand{\CompSlopeEq}{\ensuremath{-0.93}}
\newcommand{\ConvCodeRight}{\ensuremath{6.5\times10^{-19}}}
\newcommand{\ConvShiftIrregular}{\ensuremath{3.7\times10^{-3}}}
\newcommand{\ConvShiftUniform}{\ensuremath{8.7\times10^{-19}}}
\newcommand{\EoneBilinSlope}{\ensuremath{-2.02}}
\newcommand{\EoneComplexBilin}{\ensuremath{0.24}}
\newcommand{\EoneEulerMeight}{\ensuremath{0.025}}
\newcommand{\EoneEulerMone}{\ensuremath{0.21}}
\newcommand{\EoneEulerSlope}{\ensuremath{-1.03}}
\newcommand{\EoneNodtMeight}{\ensuremath{1.3}}
\newcommand{\EoneNodtMtwo}{\ensuremath{0.48}}
\newcommand{\EoneRKmax}{\ensuremath{1.6\times10^{-13}}}
\newcommand{\EoneRealBilin}{\ensuremath{9.2\times10^{-3}}}
\newcommand{\EoneZohMax}{\ensuremath{6.3\times10^{-15}}}
\newcommand{\EthreeExactChange}{\ensuremath{1.5\times10^{-16}}}
\newcommand{\EthreeExactErrEight}{\ensuremath{0.032}}
\newcommand{\EthreeMeanChange}{\ensuremath{0.30}}
\newcommand{\EthreeMeanErrEight}{\ensuremath{0.28}}
\newcommand{\EthreeMeanErrOne}{\ensuremath{0.060}}
\newcommand{\EthreeRMEneg}{\ensuremath{21}}
\newcommand{\EthreeRMEslope}{\ensuremath{-0.076}}
\newcommand{\EthreeRiemannChange}{\ensuremath{0.026}}
\newcommand{\EthreeSumChange}{\ensuremath{1.5}}
\newcommand{\EthreeTrapzChange}{\ensuremath{0.011}}
\newcommand{\EtwoCrossMax}{\ensuremath{+0.42}}
\newcommand{\EtwoCrossMin}{\ensuremath{-1.37}}
\newcommand{\EtwoEulerArtFrac}{\ensuremath{0.53}}
\newcommand{\EtwoEulerCrossFrac}{\ensuremath{+0.27}}
\newcommand{\EtwoHoldErrEight}{\ensuremath{0.030}}
\newcommand{\EtwoHoldErrOne}{\ensuremath{0.23}}
\newcommand{\EtwoNodtArtFrac}{\ensuremath{1.22}}
\newcommand{\EtwoNodtCrossFrac}{\ensuremath{-0.25}}
\newcommand{\EtwoSlope}{\ensuremath{-0.98}}
\newcommand{\EtwoZohArt}{\ensuremath{1.7\times10^{-13}}}
\newcommand{\FacAll}{\ensuremath{1.5\times10^{-16}}}
\newcommand{\FacDt}{\ensuremath{0.30}}
\newcommand{\FacDtTime}{\ensuremath{0.049}}
\newcommand{\FacDtZoh}{\ensuremath{0.30}}
\newcommand{\FacDtZohTime}{\ensuremath{0.011}}
\newcommand{\FacNone}{\ensuremath{0.46}}
\newcommand{\FacTime}{\ensuremath{0.25}}
\newcommand{\FacZoh}{\ensuremath{0.48}}
\newcommand{\FirstRunSeconds}{\ensuremath{20.6}}
\newcommand{\HfourEulerFrac}{\ensuremath{31}}
\newcommand{\HsevenOrigSlope}{\ensuremath{+0.42}}
\newcommand{\NfourFloor}{\ensuremath{1.1\times10^{-4}}}
\newcommand{\NoneNaive}{\ensuremath{8.0\times10^{-4}}}
\newcommand{\NthreeEulerLarge}{\ensuremath{4.6}}
\newcommand{\NthreeEulerSmall}{\ensuremath{3.5\times10^{-3}}}
\newcommand{\NtwoAbar}{\ensuremath{-0.977}}
\newcommand{\NtwoBilArt}{\ensuremath{1.1}}
\newcommand{\NtwoEulerArt}{\ensuremath{170}}
\newcommand{\NunitsPrimary}{\ensuremath{32}}
\newcommand{\RealDevMSE}{\ensuremath{0.012}}
\newcommand{\RealDevMax}{\ensuremath{0.84}}
\newcommand{\RealDiscSeight}{\ensuremath{0.011}}
\newcommand{\RealDiscSeightCIhi}{\ensuremath{0.011}}
\newcommand{\RealDiscSeightCIlo}{\ensuremath{0.010}}
\newcommand{\RealDiscStwo}{\ensuremath{6.1\times10^{-3}}}
\newcommand{\RealEvalSec}{\ensuremath{24}}
\newcommand{\RealJoneMSE}{\ensuremath{1.1}}
\newcommand{\RealLabelSD}{\ensuremath{2.35}}
\newcommand{\RealLayerOneDp}{\ensuremath{1.8\times10^{-14}}}
\newcommand{\RealLayerOneFp}{\ensuremath{7.2\times10^{-6}}}
\newcommand{\RealLayerTwoMax}{\ensuremath{7.2\times10^{-3}}}
\newcommand{\RealLayerTwoMin}{\ensuremath{4.1\times10^{-3}}}
\newcommand{\RealMSEczero}{\ensuremath{0.0167}}
\newcommand{\RealMSEheight}{\ensuremath{0.0156}}
\newcommand{\RealMSEseight}{\ensuremath{0.0160}}
\newcommand{\RealPoolCzero}{\ensuremath{0.063}}
\newcommand{\RealPoolLabelAbs}{\ensuremath{1.59}}
\newcommand{\RealPoolMeanHeight}{\ensuremath{0.56}}
\newcommand{\RealPoolResHeight}{\ensuremath{0.063}}
\newcommand{\RealPoolTimeHeight}{\ensuremath{0.068}}
\newcommand{\RealRelDiffHeight}{\ensuremath{-6.5}}
\newcommand{\RealRelDiffSeight}{\ensuremath{-4.1}}
\newcommand{\RealRelDiffStwo}{\ensuremath{-3.0}}
\newcommand{\RealRunNatCzero}{\ensuremath{0.053}}
\newcommand{\RealRunNatSeight}{\ensuremath{1.07}}
\newcommand{\RealRunResSeight}{\ensuremath{0.054}}
\newcommand{\RealStep}{\ensuremath{1500}}
\newcommand{\RealTrainSec}{\ensuremath{97}}


\icmltitlerunning{Where Sampling-Grid Dependence Enters Selective SSMs}

\begin{document}

\twocolumn[
  \icmltitle{Where Sampling-Grid Dependence Enters Selective State-Space Models}

  \begin{icmlauthorlist}
    \icmlauthor{Anonymous Authors}{anon}
  \end{icmlauthorlist}

  \icmlaffiliation{anon}{Anonymous Institution, Anonymous City, Anonymous Country}

  \icmlcorrespondingauthor{Anonymous Authors}{anonymous@example.com}

  \icmlkeywords{state-space models, selective SSM, discretization, irregular sampling, zero-order hold}

  \vskip 0.3in
]

\printAffiliationsAndNotice{}

\begin{abstract}
A sequence model can change its output when only the time grid of the same physical input changes.
We ask where this dependence enters a selective state-space model (SSM) and which of it is an artifact rather than new information.
In fixed-parameter toys we separate interval splitting, which keeps the held input path, from new observations, and find four entry points: a step that counts samples, an inexact input-matrix discretization, sample-count readouts, and the resampling of internal signals between layers; a physical step, exact zero-order hold and an exact time integral remove the single-layer artifact only together.
We then train the official implementation of a selective SSM whose step is already physical (TIDES) once on a small synthetic regression task and evaluate the same checkpoint on the same test trajectories: splitting the held input eight-fold changes its outputs by $\RealDiscSeight$ (relative $\ell_2$) while resampling to the training grid changes nothing by construction, the change is absent from the first block ($\RealLayerOneDp$ in float64) and appears in the second, and under a density change a sample-count pooled readout errs by $\RealPoolMeanHeight$ where time-weighted pooling of the same outputs errs by $\RealPoolTimeHeight$.
These are single-seed pilot measurements on synthetic data; we propose no new architecture.
\end{abstract}

%=====================================================================================
\section{Introduction}
\label{sec:intro}

Continuous-time state-space layers are applied to sequences by discretizing a linear ODE over each interval between observations.
When the step is the physical elapsed time, the recurrence inherits the semantics of the continuous system, which lets S4, LSSL and S5 change the sampling rate at test time or consume irregular timestamps~\citep{gu2021s4,gu2021lssl,smith2022s5}.
Selective SSMs such as Mamba compute the step from the current token~\citep{gu2023mamba,dao2024mamba2}, so it acts as a content gate; TIDES restores a physical step and moves selectivity to the state matrix~\citep{soydan2026tides}.

We do not propose another architecture.
We ask a measurement question: \emph{when only the time grid of the same physical input path changes, where does the output change enter, and which of it is an artifact?}
This requires separating two operations.
\emph{Splitting} a held interval into $m$ sub-steps that repeat the held value leaves the physical input path unchanged, so any output change is an artifact.
\emph{Adding observations} of the same signal changes the reconstructed path, so part of the change is genuine information.

\paragraph{Known and not claimed.}
The physical $\Delta t$ in the discretization, exact zero-order-hold (ZOH) formulas and flow composition are standard~\citep{gu2021lssl,smith2022s5}; that the released Mamba code uses $\bar B=\Delta B$ (``exponential-Euler'') is documented by \citet{lahoti2026mamba3}; that a continuous-time reading of selective SSMs needs $\Delta\propto dt$ is shown by \citet{cirone2024theoretical}; and keeping $\Delta$ physical with selectivity on the decay rate is the design of \citet{soydan2026tides}.

\paragraph{Contributions.}
None is an architecture.
\begin{enumerate}
  \item A protocol that separates splitting from new observations, with an exact path/artifact decomposition and explicit denominators (\cref{sec:setup}).
  \item Controlled toy measurements of four entry points and of the interaction of fixes: each fix alone leaves at least half of the change (\cref{sec:toy}).
  \item The inter-layer mechanism stated with explicit assumptions: with a physical step and exact ZOH in every layer, a stack that holds internal endpoint values still changes under splitting, at first order in the step, without new external information (\cref{sec:coupling}).
  \item A check that the TIDES paper equations and the released code use different interval pairings (\cref{sec:convention}).
  \item A single-seed pilot on a trained official TIDES checkpoint: native-grid vs.\ training-grid-resampled evaluation and a pooling-only ablation on the same test trajectories, with discrepancy, task loss and runtime, and an exploratory localization of the change to the second block (\cref{sec:real}).
\end{enumerate}

%=====================================================================================
\section{Related Work}
\label{sec:related}

\paragraph{Discretization in structured SSMs.}
S4 uses the bilinear transform and reports a test-time resolution change obtained by changing its internal step~\citep{gu2021s4}; LSSL handles irregular data by discretizing with a different timescale~\citep{gu2021lssl}; S4D gives bilinear and ZOH rules~\citep{gu2022s4d}.
S5 uses ZOH, $\bar\Lambda=e^{\Lambda\Delta}$, $\bar B=\Lambda^{-1}(\bar\Lambda-I)\tilde B$, rescales $\Delta$ globally for zero-shot 8\,kHz evaluation of a 16\,kHz model, and supplies a per-step $\Delta_t$ for an irregular pendulum task, with ablations that drop the interval or append it as a feature~\citep{smith2022s5}.
S4ND notes that only the relative rate matters~\citep{nguyen2022s4nd}; event-camera SSMs evaluate at different frequencies~\citep{zubic2024event}.
These are LTI layers.

\paragraph{Selective SSMs.}
Mamba makes $\Delta,B,C$ input-dependent and states the ZOH rule~\citep{gu2023mamba}; Mamba-2 keeps its discretization~\citep{dao2024mamba2}.
Mamba-3 derives right-endpoint ZOH for time-varying systems, names the released rule exponential-Euler, and proposes an exponential-trapezoidal rule that is second order only if its data-dependent weight is $\tfrac12+O(\Delta_t)$~\citep{lahoti2026mamba3}; its evaluation concerns language modeling and related tasks.
Liquid-S4 makes the transition input-dependent~\citep{hasani2022liquid}.

\paragraph{Physical time in selective SSMs.}
TIDES keeps $\tilde\Delta\equiv\Delta$, makes $\mathrm{Re}\,\Lambda$, $B$ and $C$ input-dependent, and uses $\bar B_k=\Lambda_k^{-1}(\bar\Lambda_k-I)B_k$~\citep{soydan2026tides}.
Its Fading Flash diagnostic rescales a global step for fixed-length sequences and shows that a Mamba surrogate with $\Delta$ as an input channel fails outside the training range; its random-drop study removes observations and reports accuracy.
In the sections we read it neither refines a fixed held path nor decomposes output changes.
It is the closest prior work; we evaluate its official code, not a re-implementation, and we treat its architecture as prior art.

\paragraph{Continuous-time models.}
Neural CDEs are driven by an interpolation of the observations~\citep{kidger2020ncde}; Log-NCDEs build on NCDEs, which are robust to irregular sampling~\citep{walker2024logncde}; \citet{cirone2024theoretical} read selective SSMs as linear CDEs in a limit where $\delta$ plays the role of $dt$; oscillatory SSMs study implicit and IMEX integration~\citep{rusch2024linoss}.
We did not find, in these sources, a refinement of the same held path for selective SSMs that separates splitting from new observations and treats readouts and inter-layer resampling as entry points; this is absence of evidence, not proof of novelty.

%=====================================================================================
\section{Problem Setup}
\label{sec:setup}

\paragraph{Model and holds.}
The toy is a scalar-input selective model with $N$ diagonal modes,
\begin{equation}
\dot h_n=g(u)\big(\lambda_n h_n+B_n(u)\,u\big),\qquad y=\mathrm{Re}\textstyle\sum_nC_n(u)h_n,
\label{eq:ct}
\end{equation}
with $\mathrm{Re}\,\lambda_n<0$, a gate $g(u)=\mathrm{softplus}(wu+\beta)$ shared across modes, and affine $B(u),C(u)$.
A grid $0=t_0<\dots<t_K=T$ has steps $dt_k$ and mesh $\delta$.
We use the right-endpoint hold $\tilde u(t)=u(t_k)$ on $(t_{k-1},t_k]$, which is the convention of the Mamba recurrence and of the TIDES code (\cref{sec:convention}).
A discrete variant is $h_k=\bar A_kh_{k-1}+\bar B_ku_k$ with $\Delta_k=dt_k\,g(u_k)$ (``$\Delta{=}dt\,g$'') or $\Delta_k=\tau g(u_k)$ with a fixed $\tau$ (``$\Delta{=}\tau g$''), and exact ZOH, Euler-$B$ ($\bar B=\Delta B$) or bilinear rules; $\tau$ is the base step.

\paragraph{Grid changes and decomposition.}
\emph{Splitting} divides base intervals into $m$ equal sub-steps that all receive the held value; \emph{new observations} give each new sub-time the value of the continuous signal.
Outputs are compared at the common base times.
With $D_m$ the outputs of a variant on grid $m$ and $Y_m$ the continuous-time solution of \eqref{eq:ct} on the held path of grid $m$, the change splits exactly as
\begin{equation}
T_m=\underbrace{(Y_m-Y_1)}_{P_m\text{: path}}+\underbrace{\big[(D_m-Y_m)-(D_1-Y_1)\big]}_{R_m\text{: artifact}},
\label{eq:decomp}
\end{equation}
so $\|T_m\|^2=\|P_m\|^2+\|R_m\|^2+2\langle P_m,R_m\rangle$.
Under splitting $P_m=0$.
Because the cross term can be large and of either sign, we never report a ratio of norms as a share.

\paragraph{Readouts.}
We distinguish sample-count readouts ($\sum_ky_k$, $\frac1K\sum_ky_k$), time-weighted endpoint sums (right-endpoint $\frac1T\sum_ky_k\,dt_k$ and trapezoidal), and the exact held-path integral $\frac1T\int_0^Ty\,dt$, defined only for exact-ZOH states.

%=====================================================================================
\section{Where Grid Dependence Enters: Analysis}
\label{sec:analysis}

\subsection{Single layer}
\label{sec:single}

\begin{proposition}[standard]
\label{prop:split}
Fix a base interval $(t_{k-1},t_k]$ and suppose
(i) every sub-step receives the held value $u_k$;
(ii) every coefficient ($g,B,C,\lambda$) is a function of the held value only;
(iii) the step is (sub-step length)$\,\cdot g(u_k)$; and
(iv) $\bar A,\bar B$ are the exact ZOH formulas.
Then, from the same state at $t_{k-1}$, the state and output at $t_k$ are the same for every $m$; by induction the outputs at all base times are split-invariant.
\end{proposition}
\begin{proof}
Each sub-step is the exact flow $\Phi_s$ of $\dot h=g(u_k)(\lambda h+B(u_k)u_k)$ over its length $s$, and $\Phi_{s_1}\circ\Phi_{s_2}=\Phi_{s_1+s_2}$.
\end{proof}

The proposition is textbook; its use is that each entry point violates one assumption.
$\Delta{=}\tau g$ violates (iii): $m$ sub-steps apply $m$ times the decay.
Euler-$B$ violates (iv), with local error $|\Delta Bu|\,|\varphi_1(\Delta\lambda)-1|\le|Bu|\Delta^2|\lambda|/2$, $\varphi_1(z)=(e^z-1)/z$.
(ii) fails whenever a coefficient sees anything besides the held value: a token-indexed convolution (Mamba's \texttt{d\_conv}${=}4$), time or position features, state-dependent transitions, or training-mode normalization statistics over time.
Bilinear is stable for $\mathrm{Re}\,z<0$ at every step, but for real $z<-2$ its transition is negative (sign-alternating transients) and it tends to $-1$ instead of $e^z\to0$ (inaccurate stiff decay); its held-input fixed point is exact.

\subsection{Several layers: internal resampling}
\label{sec:coupling}

In a stack, layer~2 is driven by the continuous output of layer~1, but a layerwise recurrence gives it only endpoint samples held over each step.
Splitting the input grid changes which internal samples it sees, without adding any external observation.
The minimal instance is
\begin{equation}
\dot h_1=-a h_1+b\,u(t),\qquad \dot h_2=-c\,h_2+d\,h_1(t).
\label{eq:cascade}
\end{equation}

\begin{proposition}[standard]
\label{prop:cascade}
Let $u$ be held on a grid with mesh $\delta$, $|u|\le U$, $a\ge0$, $c>0$ and $h_1(0)=h_2(0)=0$.
(a) The exact coupled solution of \eqref{eq:cascade} at base times is split-invariant.
(b) The layerwise computation that is exact in layer~1 and holds $h_1(t_k)$ on $(t_{k-1},t_k]$ in layer~2 has absolute error $|e_k|\le|b\,d|\,U\,\delta\,e^{c\delta}/c$ at every grid time, for every $T$ and also for $a=c$.
\end{proposition}
\begin{proof}
(a) follows as in \cref{prop:split}, since the coupled system has constant coefficients on each held interval.
(b) Layer~1 is exact, so with $\tau_k=dt_k$,
$e_k=e^{-c\tau_k}e_{k-1}+d\int_0^{\tau_k}e^{-c(\tau_k-s)}[h_1(t_k)-h_1(t_{k-1}+s)]\,ds$.
If $a>0$ then $|h_1|\le|b|U/a$, so $a|h_1|\le|b|U$; if $a=0$ then $a|h_1|=0$; either way $|\dot h_1|\le2|b|U$.
On each held interval $h_1$ is $C^1$, so the bracket is at most $2|b|U(\tau_k-s)$, and with $e^{-c(\tau_k-s)}\le1$ the local term is at most $|b\,d|U\tau_k^2$.
Unrolling from $e_0=0$ gives $|e_k|\le|b\,d|U\delta\sum_{j\le k}\tau_je^{-c(t_k-t_j)}$, and since $e^{-c(t_k-t_j)}\le e^{c\delta}e^{-c(t_k-s)}$ for $s\in(t_{j-1},t_j]$, the sum is at most $e^{c\delta}\int_0^{t_k}e^{-c(t_k-s)}ds\le e^{c\delta}/c$.
\end{proof}

\begin{remark}[scope of \cref{prop:cascade}]
\label{rem:scope}
The argument is elementary and has not been checked by a third party.
A nonzero $h_1(0)$ adds $a|h_1(0)|$ to $2|b|U$; a nonzero $h_2(0)$ does not enter.
Uniformity in $T$ needs $c>0$: as $c\to0$ the sum is bounded only by $t_k$ and the bound grows linearly in time; $a<0$ (an unstable first layer) is not covered.
In the stored P1-COMP-01 outputs, the largest error at the base times is 9--20\% of the bound for $m\in\{1,2,4,8\}$, which is consistent with, not a proof of, the bound.
The statement is for a linear time-invariant scalar cascade; it gives no guarantee for input-dependent coefficients, nonlinear maps between layers, or trained deep models, where the Lipschitz and contraction constants it uses are not controlled.
\end{remark}

When all conditions of \cref{prop:split} hold in every layer, a stacked recurrence on the observation grid can therefore still change under splitting because the inter-layer signal is resampled; this is a reconstruction error of an \emph{internal} signal, not new information.
It vanishes if the inter-layer signal is represented exactly or if the stack runs on a grid that does not depend on the observations.

\subsection{Readouts}
The exact held-path integral is additive over sub-intervals and hence split-invariant under exact-ZOH states; endpoint sums converge to it only where the grid is refined; sample-count readouts weight samples, not time, so the sample mean converges to a density-weighted average.

\subsection{Interval conventions in the TIDES paper and code}
\label{sec:convention}

The TIDES paper writes $x_{k+1}=\bar\Lambda(\Delta_k)x_k+\bar B(\Delta_k)u_k$ with $\Delta_k=t_{k+1}-t_k$ and $y_k=Cx_k+Du_k$, i.e.\ $u_k$ is held forward on $[t_k,t_{k+1})$ and the state in $y_k$ does not yet contain $u_k$~\citep{soydan2026tides}.
The pinned code scans $z_k=\bar\Lambda(t_k-t_{k-1})z_{k-1}+\bar B(t_k-t_{k-1})u_k$ with $y_k=\mathrm{Re}(Cz_k)+Du_k$, and its forecasting collate builds the step sequence $[t_0,t_1-t_0,\dots]$: $u_k$ is held backward on $(t_{k-1},t_k]$ and the state in $y_k$ contains it.
Both are exact ZOH solutions, but of different input reconstructions.
On a uniform grid every gap is equal, so $z_k=x_{k+1}$: the difference is a one-index shift of the state, which changes the alignment between the state part of $y_k$ and the feedthrough $Du_k$.
On an irregular grid each value is paired with a different gap, so the difference is not a relabeling.
With the pinned code, untrained weights and a fixed impulse input, the code equals the backward-hold recurrence to $\ConvCodeRight$; the one-index shift reproduces the paper form to $\ConvShiftUniform$ on a uniform grid but leaves $\ConvShiftIrregular$ on an irregular one.
The released code, not the printed equations, defines the model behind TIDES' reported results, and it coincides with the right-endpoint convention used throughout this paper; our splitting protocol repeats the right-endpoint value accordingly.
We make no claim that this affects TIDES' results.

%=====================================================================================
\section{Trained Official TIDES: a Single-Seed Pilot}
\label{sec:real}

\paragraph{Protocol (fixed before training; \cref{app:real}).}
\emph{Model:} official TIDES at a pinned commit (MIT), $\AdapterParams$ parameters, two blocks, input-dependent $\mathrm{Re}\,\Lambda$, $B$, $C$, exact ZOH, causal, no convolution, a linear per-step head; CPU, float32, 2 threads.
\emph{Task:} regression of the teacher \eqref{eq:ct} with fixed parameters, driven by sine-sum inputs, over $T{=}8$ on a uniform training grid of $K{=}32$ steps; labels are the teacher output on the continuous path at the 32 grid times (label s.d.\ $\RealLabelSD$).
\emph{Split:} 896 trajectory identifiers split into 512/128/256 train/calibration/test \emph{before} any grid is built; normalization statistics from the training split only.
\emph{Training:} one seed (0), Adam, 2000 steps, batch 32; calibration MSE every 100 steps; the lowest-calibration-MSE state (step $\RealStep$, MSE $\RealDevMSE$) was saved, hashed and frozen before any test trajectory was generated ($\RealTrainSec$\,s in total).
\emph{Conditions:} C0 (training grid); S2/S4/S8, which split every interval with the held value repeated; H8, which splits only the first half by 8 (a density change on the same held path); and, as a secondary non-refinement condition, J1 with $K$ random observation times.
\emph{Arms on the same checkpoint and test trajectories:} \emph{native} (the condition's grid with physical step sizes), \emph{resampled} (last observation carried forward onto the training grid, which uses only observed values and never the continuous path), and, for the pooled readout, \emph{pooling-only} (the native outputs pooled by the time-weighted sum instead of the sample mean).
\emph{Equivalence rule:} the unit is the per-trajectory test MSE in label units; the margin is $\pm5\%$ of the resampled arm's mean MSE for the same condition, applied to the mean paired difference; ``no practical difference'' requires the 95\% bootstrap CI (2000 resamples of test trajectories) to lie inside the margin.
The 5\% value is a convention fixed before any trained result (about 2.5\% in RMSE), not derived from an application cost; non-significance is never read as equivalence.

% AUTO-GENERATED by scripts/make_paper_assets.py from results/raw -- do not edit by hand.
\begin{table}[t]
\caption{DEVELOPMENT run (P1-REAL-01): trained official TIDES, seed 0, on its own test set (256 trajectories, ids 640--895; one checkpoint selected on the calibration split; paired).
``disc.'': mean relative $\ell_2$ change of the per-step outputs at the 32 training-grid times against the same arm on C0.
MSE: per-trajectory mean squared error to the teacher at those times (original units; label s.d.\ 2.35).
$\Delta$MSE: mean paired difference native $-$ resampled in units of $10^{-4}$ (label units squared) with a 95\% bootstrap CI; the margin is $\pm5\%$ of the point-estimated resampled mean MSE of the same condition (exploratory rule of the development run).
Verdict (native vs.\ resampled MSE): lower/higher if the CI excludes 0 on that side. Runtimes of this run are not tabulated; they were forward-only and are superseded by \cref{tab:timing}.
C0 is the training grid; S$m$ and H8 (first half split by 8) keep the held path; $^\dagger$J1 (random observation times) is a secondary non-refinement condition whose held path differs by construction.}
\label{tab:real}
\centering\small
\begin{tabular}{lcccccc}
\toprule
cond. & disc.\ native & disc.\ resampled & MSE native & MSE resampled & $\Delta$MSE [95\% CI] ($10^{-4}$) & margin rule \\
\midrule
C0 & $0$ & $0$ & $0.0167$ & $0.0167$ & $0$ $[0,\,0]$ & identical inputs (check) \\
S2 & $6.1\times10^{-3}$ & $0$ & $0.0162$ & $0.0167$ & $-5.1$ $[-7.3,\,-2.7]$ & within $\pm5\%$ \\
S4 & $9.2\times10^{-3}$ & $0$ & $0.0161$ & $0.0167$ & $-6.5$ $[-9.9,\,-2.9]$ & lower \\
S8 & $0.011$ & $0$ & $0.0160$ & $0.0167$ & $-6.9$ $[-10.8,\,-2.7]$ & lower \\
H8 & $7.8\times10^{-3}$ & $0$ & $0.0156$ & $0.0167$ & $-10.8$ $[-14.0,\,-7.6]$ & lower \\
J1$^\dagger$ & $0.35$ & $0.39$ & $1.14$ & $1.22$ & $-785$ $[-2120,\,+623]$ & inconclusive \\
\bottomrule
\end{tabular}
\end{table}


\paragraph{Discrepancy, task loss and runtime.}
\Cref{tab:real} reports all three together.
Resampling makes the model split-invariant by construction: its outputs on S2--S8 and H8 equal those on C0 exactly.
Native evaluation is not: splitting every interval changes the outputs by $\RealDiscStwo$ at $m{=}2$ and $\RealDiscSeight$ at $m{=}8$ (95\% CI $[\RealDiscSeightCIlo,\RealDiscSeightCIhi]$), far above the float32 floor of about $10^{-4}$ (\cref{app:numerics}).
The native task loss on the refined grids is slightly \emph{lower} than on the training grid (\RealRelDiffStwo\% at S2, \RealRelDiffSeight\% at S8 and \RealRelDiffHeight\% at H8 relative to the resampled MSE of $\RealMSEczero$).
By the pre-specified rule, S2 shows no practical difference; for S4, S8 and H8 the interval excludes zero but does not lie beyond the margin, so neither equivalence nor a practical difference is established.
We do not know why the refined native computation is slightly more accurate here; a finer internal sampling of the inter-layer signal is one untested explanation.
Native inference cost grows with the number of tokens (from $\RealRunNatCzero$\,s to $\RealRunNatSeight$\,s per batch from C0 to S8 in this implementation), while resampling stays at about $\RealRunResSeight$\,s.
On the secondary J1 condition, the errors of both arms are much larger (native MSE $\RealJoneMSE$ against $\RealMSEczero$ on C0) because $K$ random times observe the input differently, and the paired difference is inconclusive.

% AUTO-GENERATED by scripts/make_paper_assets.py from results/raw -- do not edit by hand.
\begin{table}[t]
\caption{Pooled readout of the same trained checkpoint (P1-REAL-01 pilot). Label: exact time average of the teacher output over $[0,T]$ (mean absolute value 1.59).
``mean'': sample-count mean over tokens (as in the TIDES classifier head); ``time'': right-endpoint time-weighted sum over the same native outputs (pooling-only ablation).
Mean absolute error to the label and mean absolute change against C0, over 256 test trajectories.}
\label{tab:pool}
\centering\small
\resizebox{\columnwidth}{!}{%
\begin{tabular}{lcccccc}
\toprule
 & \multicolumn{3}{c}{error to pooled label} & \multicolumn{3}{c}{change vs.\ C0} \\
\cmidrule(lr){2-4}\cmidrule(lr){5-7}
cond. & native mean & native time & resampled & native mean & native time & resampled \\
\midrule
C0 & $0.063$ & $0.063$ & $0.063$ & $0$ & $0$ & $0$ \\
S8 & $0.070$ & $0.070$ & $0.063$ & $0.032$ & $0.032$ & $0$ \\
H8 & $0.56$ & $0.068$ & $0.063$ & $0.55$ & $0.016$ & $0$ \\
J1 & $0.37$ & $0.28$ & $0.29$ & $0.36$ & $0.26$ & $0.30$ \\
\bottomrule
\end{tabular}}
\end{table}


\paragraph{Pooled readout.}
\Cref{tab:pool} evaluates a sequence-level readout of the same checkpoint: the time average of the teacher output (mean magnitude $\RealPoolLabelAbs$).
On uniform splits, sample-mean and time-weighted pooling coincide.
Under the density change H8, the sample mean over native tokens (the pooling used by the TIDES classifier head) errs by $\RealPoolMeanHeight$, whereas time-weighted pooling of the \emph{same} native outputs errs by $\RealPoolTimeHeight$, close to the resampled arm ($\RealPoolResHeight$) and to C0 ($\RealPoolCzero$).
In this pilot the pooled-readout grid dependence is therefore almost entirely a readout effect, which the pooling-only change removes; the pre-registered direction (R-P5) was supported.

% AUTO-GENERATED by scripts/make_paper_assets.py from results/raw -- do not edit by hand.
\begin{table}[t]
\caption{Where the splitting change appears in the development checkpoint (EXPLORATORY, added after its results were seen). Mean relative $\ell_2$ change against C0 at the training-grid times of the
features after each block and of the (normalized) output, with the frozen weights evaluated in float32 and cast to float64. The encoder output is unchanged (0) in all cases.}
\label{tab:layers}
\centering\small
\begin{tabular}{lcccccc}
\toprule
 & \multicolumn{2}{c}{block 1} & \multicolumn{2}{c}{block 2} & \multicolumn{2}{c}{output} \\
\cmidrule(lr){2-3}\cmidrule(lr){4-5}\cmidrule(lr){6-7}
cond. & fp32 & fp64 & fp32 & fp64 & fp32 & fp64 \\
\midrule
S2 & $5.6\times10^{-6}$ & $5.1\times10^{-15}$ & $4.1\times10^{-3}$ & $4.1\times10^{-3}$ & $6.8\times10^{-3}$ & $6.8\times10^{-3}$ \\
S4 & $6.0\times10^{-6}$ & $1.1\times10^{-14}$ & $6.1\times10^{-3}$ & $6.1\times10^{-3}$ & $0.010$ & $0.010$ \\
S8 & $7.2\times10^{-6}$ & $1.8\times10^{-14}$ & $7.2\times10^{-3}$ & $7.2\times10^{-3}$ & $0.012$ & $0.012$ \\
H8 & $5.8\times10^{-6}$ & $1.4\times10^{-14}$ & $5.5\times10^{-3}$ & $5.4\times10^{-3}$ & $9.1\times10^{-3}$ & $9.1\times10^{-3}$ \\
\bottomrule
\end{tabular}
\end{table}


\paragraph{Where the change appears (exploratory).}
After seeing the results we measured the change at intermediate features (\cref{tab:layers}); this analysis was not pre-registered.
The encoder output is unchanged, the first block changes by at most $\RealLayerOneFp$ in float32 and $\RealLayerOneDp$ with the weights cast to float64, and the second block changes by $\RealLayerTwoMin$ to $\RealLayerTwoMax$ in both precisions.
This is the pattern that \cref{prop:split,prop:cascade} predict: a first layer that satisfies the single-layer conditions is split-invariant up to rounding, and the change enters when the second block receives additional internal samples.
It localizes the effect in one checkpoint; it does not prove the mechanism for trained stacks in general.

\paragraph{What the pilot does not show.}
One seed, one small model and one synthetic task are a conditional pilot, not evidence about the performance distribution of the training algorithm or about TIDES on its benchmarks; the teacher belongs to the same continuous-time family as the model, which may favour it; and the calibration curve was unstable (calibration MSE up to $\RealDevMax$ between evaluations), so the selected checkpoint may not be representative.
Nothing here transfers to real sensor or clinical streams.

%=====================================================================================
\section{Controlled Toy Results}
\label{sec:toy}

These fixed-parameter results are numerical checks and controlled measurements, not proofs, and say nothing directly about trained models.
The pre-registered FIRST\_RUN used \NunitsPrimary{} independent units (one parameter and input draw each), $T{=}8$, $K{=}16$, $m\in\{1,2,4,8\}$, 2 CPU threads and \FirstRunSeconds\,s; details and deviations are in \cref{app:prereg}.

\paragraph{Splitting.}
With $\Delta{=}dt\,g$ and exact ZOH the output is split-invariant to $\EoneZohMax$ and agrees with an independent RK4 reference to $\EoneRKmax$ (an engineering check of \cref{prop:split}).
Euler-$B$ with the correct step converges at first order (median slope $\EoneEulerSlope$) but has a median base-step artifact of $\EoneEulerMone$ in this toy; the magnitude depends on the toy's $\Delta|\lambda|$ scale.
Bilinear converges at second order ($\EoneBilinSlope$); with $\Delta{=}\tau g$ the artifact grows from $\EoneNodtMtwo$ to $\EoneNodtMeight$ (medians, $m{=}2\to8$), by construction (\cref{tab:split}, \cref{fig:convergence}a).
In a stiff mode ($\lambda{=}-500$) the bilinear transition is $\NtwoAbar$ and Euler-$B$ has an artifact of $\NtwoEulerArt$ (\cref{app:toy}).

% AUTO-GENERATED by scripts/make_paper_assets.py from results/raw -- do not edit by hand.
\begin{table}[t]
\caption{Splitting a held interval into $m$ sub-steps (FR-E1-SPLIT; toy, 32 units, uniform base grid, real modes).
Entries are the median over units of the output artifact $\|y_m-y^\star\|_2/\|y^\star\|_2$ at the 17 common base times, where $y^\star$ is the exact held-path solution;
slope is the median per-unit log--log slope over $m\in\{1,2,4,8\}$. $\tau$ is the base step, so $\Delta{=}\tau g$ coincides with $\Delta{=}dt\,g$ at $m{=}1$.}
\label{tab:split}
\centering\small
\begin{tabular}{lcccccc}
\toprule
variant & $m{=}1$ & $m{=}2$ & $m{=}4$ & $m{=}8$ & max, $m{=}8$ & slope \\
\midrule
exact ZOH, $\Delta{=}dt\,g$ & $0$ & $2.6\times10^{-16}$ & $3.8\times10^{-16}$ & $5.5\times10^{-16}$ & $6.3\times10^{-15}$ & -- \\
Euler-$B$, $\Delta{=}dt\,g$ & $0.21$ & $0.10$ & $0.050$ & $0.025$ & $0.12$ & $-1.03$ \\
bilinear, $\Delta{=}dt\,g$ & $9.5\times10^{-3}$ & $2.3\times10^{-3}$ & $5.6\times10^{-4}$ & $1.4\times10^{-4}$ & $8.1\times10^{-4}$ & $-2.02$ \\
exact ZOH, $\Delta{=}\tau g$ & $0$ & $0.48$ & $0.93$ & $1.3$ & $2.8$ & -- \\
Euler-$B$, $\Delta{=}\tau g$ & $0.21$ & $0.68$ & $1.1$ & $1.4$ & $3.1$ & $0.93$ \\
bilinear, $\Delta{=}\tau g$ & $9.5\times10^{-3}$ & $0.48$ & $0.94$ & $1.3$ & $2.8$ & $2.25$ \\
\bottomrule
\end{tabular}
\end{table}

% AUTO-GENERATED by scripts/make_paper_assets.py from results/raw -- do not edit by hand.
\begin{figure}[t]
\centering
\definecolor{seriesA}{HTML}{2A78D6}\definecolor{seriesB}{HTML}{EB6834}\definecolor{seriesC}{HTML}{1BAF7A}
\begin{tikzpicture}
\begin{groupplot}[group style={group size=2 by 1, horizontal sep=1.3cm},
  width=0.54\columnwidth, height=0.5\columnwidth, xmode=log, ymode=log, log basis x=2,
  xtick={1,2,4,8}, xticklabels={1,2,4,8}, xlabel={sub-steps $m$}, tick label style={font=\scriptsize},
  label style={font=\scriptsize}, title style={font=\scriptsize}, grid=major, grid style={gray!20, line width=0.3pt},
  axis line style={gray!60}, legend style={font=\tiny, draw=none, fill=none}, legend cell align=left]
\nextgroupplot[title={(a) single layer, FR-E1}, ylabel={relative error}, legend pos=south west, ymin=5e-5, ymax=5]
\addplot[seriesA, line width=1pt, mark=*, mark size=1.6pt] coordinates {(1,2.103473e-01) (2,1.022642e-01) (4,5.040279e-02) (8,2.495612e-02)};
\addlegendentry{Euler-$B$, $\Delta{=}dt\,g$}
\addplot[seriesB, line width=1pt, dashed, mark=square*, mark size=1.6pt] coordinates {(1,9.454201e-03) (2,2.284158e-03) (4,5.578572e-04) (8,1.386879e-04)};
\addlegendentry{bilinear, $\Delta{=}dt\,g$}
\addplot[seriesC, line width=1pt, dotted, mark=triangle*, mark size=2pt] coordinates {(2,4.762033e-01) (4,9.339680e-01) (8,1.275009e+00)};
\addlegendentry{exact ZOH, $\Delta{=}\tau g$}
\nextgroupplot[title={(b) two-layer cascade, P1-COMP-01}, legend pos=south west, ymin=5e-3, ymax=0.5]
\addplot[seriesA, line width=1pt, mark=*, mark size=1.6pt] coordinates {(1,1.884736e-01) (2,1.001190e-01) (4,5.162333e-02) (8,2.621494e-02)};
\addlegendentry{$a{=}1,\,c{=}0.5$}
\addplot[seriesB, line width=1pt, dashed, mark=square*, mark size=1.6pt] coordinates {(1,1.966501e-01) (2,1.063295e-01) (4,5.535558e-02) (8,2.825132e-02)};
\addlegendentry{$a{=}c{=}1$}
\end{groupplot}
\end{tikzpicture}
\caption{Error against the exact held-path solution as the same held interval is split into $m$ sub-steps (medians over 32 toy units in (a); one deterministic case per line in (b)).
Exact ZOH with $\Delta{=}dt\,g$ (a) and the exact coupled solution (b) stay at $10^{-15}$ and are not drawn. The $\Delta{=}\tau g$ line starts at $m{=}2$ because it is exact at $m{=}1$ by construction.
Values are listed in \cref{tab:split,tab:coupling}.}
\label{fig:convergence}
\end{figure}


\paragraph{New observations.}
With exact ZOH the artifact against RK4 on the new held path stays below $\EtwoZohArt$, so the whole change is the change of the continuous-time solution on the new path, i.e.\ new external information.
The exact decomposition \eqref{eq:decomp} of the other variants (exploratory, \cref{app:toy}) has cross terms between $\EtwoCrossMin$ and $\EtwoCrossMax$ of the per-unit total, so no per-unit share is defined.

\paragraph{Readouts and interaction of fixes.}
On a split-invariant backbone, a density change moves the sample mean by $\EthreeMeanChange$ and the exact integral by $\EthreeExactChange$ (scale-normalized, $m{=}8$).
The pre-registered clause H7c was mis-specified before the run and failed (slope $\HsevenOrigSlope$); its restatement passed only weakly (median slope $\EthreeRMEslope$, $\EthreeRMEneg$ of 32 units negative).
One fix at a time (\cref{tab:interaction}): exact ZOH alone gives $\FacZoh$, the physical step alone $\FacDt$, a time readout alone $\FacTime$, against $\FacNone$ without fixes; only all three together reach $\FacAll$.
Under the pre-registered stop condition this means the toy gives no support for a new architecture; the stop condition is an operational rule, not a statistical, equivalence or novelty test.

% AUTO-GENERATED by scripts/make_paper_assets.py from results/raw -- do not edit by hand.
\begin{table}[t]
\caption{One fix at a time (FR-E3-READOUT, split of the first half, $m{=}8$; median scale-normalized readout change over 32 units). The effects are not additive,
so no per-fix share is reported.}
\label{tab:interaction}
\centering\small
\begin{tabular}{lc}
\toprule
fixes applied & change \\
\midrule
none ($\Delta{=}\tau g$, Euler-$B$, sample mean) & $0.46$ \\
time readout only (trapezoidal) & $0.25$ \\
exact ZOH $B$ only & $0.48$ \\
$\Delta{=}dt\,g$ only & $0.30$ \\
$\Delta{=}dt\,g$ + trapezoidal & $0.049$ \\
$\Delta{=}dt\,g$ + exact ZOH (sample mean) & $0.30$ \\
$\Delta{=}dt\,g$ + exact ZOH + trapezoidal & $0.011$ \\
$\Delta{=}dt\,g$ + exact ZOH + exact integral & $1.5\times10^{-16}$ \\
\bottomrule
\end{tabular}
\end{table}


\paragraph{Two-layer cascade.}
For \eqref{eq:cascade} (P1-COMP-01; held input of FIRST\_RUN unit~0; \CompSeconds\,s), the exact coupled solution is split-invariant to $\CompInv$, while the layerwise computation has error $\CompErrOne$ at the base step and $\CompErrEight$ at $m{=}8$ (slope $\CompSlope$; $\CompSlopeEq$ for $a=c$; \cref{tab:coupling}, \cref{fig:convergence}b).

% AUTO-GENERATED by scripts/make_paper_assets.py from results/raw -- do not edit by hand.
\begin{table}[t]
\caption{Two-layer linear cascade $\dot h_1=-ah_1+bu$, $\dot h_2=-ch_2+dh_1$ under interval splitting (P1-COMP-01; $b{=}d{=}1$; held input of FIRST\_RUN unit 0).
Columns $m$: error of the layerwise computation (layer 2 holds the layer-1 endpoint value) against the exact coupled solution, $\|h_2^{\mathrm{lw}}-h_2\|/\|h_2\|$ over base times.
``change'' is the layerwise change $m{=}8$ vs.\ $m{=}1$; the last column is the largest change of the exact coupled solution over $m$. No external observation is added in this experiment.}
\label{tab:coupling}
\centering\small
\resizebox{\columnwidth}{!}{%
\begin{tabular}{lccccccc}
\toprule
case & $m{=}1$ & $m{=}2$ & $m{=}4$ & $m{=}8$ & slope & change & exact inv. \\
\midrule
$a{=}1,\;c{=}0.5$ & $0.19$ & $0.10$ & $0.052$ & $0.026$ & $-0.95$ & $0.15$ & $1.7\times10^{-15}$ \\
$a{=}c{=}1$ & $0.20$ & $0.11$ & $0.055$ & $0.028$ & $-0.93$ & $0.16$ & $7.1\times10^{-16}$ \\
\bottomrule
\end{tabular}}
\end{table}


%=====================================================================================
\section{Limitations}
\label{sec:limits}

\emph{Trained evidence is a pilot:} one seed, one checkpoint, one small two-block model, one synthetic task whose teacher is from the model's own family; no statement about the training algorithm, TIDES' benchmarks, Mamba, S4 or S5 is supported.
\emph{Margin:} the $\pm5\%$ equivalence margin is a convention, not an application-derived tolerance; all intervals are reported.
\emph{Entry points not measured:} token-indexed convolution, normalization in training mode, and deeper or wider stacks.
\emph{Analysis:} \cref{prop:split,prop:cascade} are standard; \cref{prop:cascade} covers only a linear time-invariant cascade (\cref{rem:scope}); the hold and Euler-$B$ bounds in \cref{app:proofs} are unproved sketches.
\emph{Exploratory parts:} the layer localization, the toy decomposition and the stiffness columns were analysed after the results were seen.
\emph{Numerics:} the float32 floor comes from an emulation.
\emph{Literature:} the key equations and experiment sections of TIDES, Mamba-3 and S5 were read in the original text; other entries were checked by keyword search.

%=====================================================================================
\section{Conclusion}
\label{sec:conclusion}

Grid dependence in selective SSMs enters through a step that counts samples, an inexact input-matrix discretization, sample-count readouts and, in stacks, the resampling of internal signals.
In controlled toys a physical step, exact ZOH and an exact integral remove the single-layer artifact, but only together, and they leave the inter-layer resampling error, which is first order in the step and carries no new information.
In a single-seed pilot on the official TIDES code, whose step is physical, the first block is split-invariant up to rounding, the change appears in the second block, resampling to the training grid removes it at the cost of discarding observations, and time-weighted pooling removes the pooled-readout effect of a density change.
Whether these magnitudes hold across seeds, tasks, sizes and implementations is open.

%=====================================================================================
\section*{Impact Statement}

This paper presents work whose goal is to advance the field of Machine Learning, specifically the understanding of how sequence models respond to changes in sampling.
All experiments use synthetic data and small CPU computations; the trained-model results come from one seed on one synthetic task.
The main risk we see is over-reading such results: a model that behaves consistently in a controlled setting may still behave differently on irregularly sampled real data, for example clinical or sensor streams, so any use in such settings needs evaluation on the actual data and sampling process.
We are not aware of other societal consequences that must be specifically highlighted here.

\bibliography{references}
\bibliographystyle{icml2026}

%=====================================================================================
\newpage
\appendix
\onecolumn

\section{Pre-registration, deviations and compute}
\label{app:prereg}

\paragraph{FIRST\_RUN.}
Hypotheses H1--H8, metrics, 32 units, 2 CPU threads, a 120\,s budget and stop conditions were committed before any experiment code ran.
Amendments recorded before any experiment ran: A1 (adaptive RK4 reference step), A2 (a restated H7c, keeping the mis-specified original), A3 (secondary configurations fixed to seeds 0--7, descriptive).
Deviations: a timing-only smoke test on seed 999 (metrics not inspected); a repeated execution after a manifest bookkeeping bug, with byte-identical data files; a partial execution killed by a shell pipe, whose outputs were deleted.
P1-COMP-01 was pre-registered separately before its code ran.
\Cref{tab:hypotheses} lists every pre-registered verdict, including the failure.

\paragraph{P1-REAL-01.}
The design, the pinned implementation and the predictions were committed before any installation.
Amendments R-A1--R-A9 were committed after an untrained adapter check and a timing smoke on four training trajectories, and before training: calibration-based checkpoint selection, same-held-path refinements only (new-observation conditions removed, one random-time condition kept as secondary), a resampler restricted to observed values, the pooled readout and pooling-only arm, the explicit equivalence unit, the single-seed pilot label, time caps, and seeding of all random generators (the TIDES initialization also draws from NumPy's global generator).
The pre-registered predictions were: native splitting discrepancy at S8 above $10^{-4}$ (supported), resampled splitting discrepancy exactly zero (supported), native runtime growing with tokens (observed), and a larger sample-mean than time-weighted pooled change on H8 (supported).

\paragraph{Compute.}
All on CPU (2 threads for experiments), no GPU, no paid API, no external data.
FIRST\_RUN: two executions of about 21\,s, one partial execution and a smoke test of about 3\,s together.
P1-COMP-01: \CompSeconds\,s.
P1-REAL-01: installation of a CPU tensor library (35\,s) and of a \LaTeX\ toolchain (66\,s); adapter check and timing smoke (${\approx}4$\,s; $\AdapterStep$\,s per batch-32 step); convention check (${\approx}2$\,s); training including label generation ($\RealTrainSec$\,s); evaluation ($\RealEvalSec$\,s); layer diagnostic (${\approx}11$\,s).

% AUTO-GENERATED by scripts/make_paper_assets.py from results/raw -- do not edit by hand.
\begin{table}[h]
\caption{Pre-registered check verdicts, including the failed H7c\_original (mis-specified before the run; amendment A2) and the weakly supported H7c\_restated. These are bookkeeping, not scientific results.}
\label{tab:hypotheses}
\centering\small
\begin{tabular}{llc}
\toprule
run & check & verdict \\
\midrule
FIRST\_RUN & H1 & PASS \\
FIRST\_RUN & H2 & PASS \\
FIRST\_RUN & H3 & PASS \\
FIRST\_RUN & H4 & PASS \\
FIRST\_RUN & H5a & PASS \\
FIRST\_RUN & H5b & PASS \\
FIRST\_RUN & H6 & PASS \\
FIRST\_RUN & H7a & PASS \\
FIRST\_RUN & H7b & PASS \\
FIRST\_RUN & H7c\_original & FAIL \\
FIRST\_RUN & H7c\_restated & PASS \\
FIRST\_RUN & H8 & PASS \\
P1-COMP-01 & COMP-H1 (distinct) & PASS \\
P1-COMP-01 & COMP-H2 (distinct) & PASS \\
P1-COMP-01 & COMP-H3 (distinct) & PASS \\
P1-COMP-01 & COMP-H4 (distinct) & PASS \\
P1-COMP-01 & COMP-H5 (distinct) & PASS \\
P1-COMP-01 & COMP-H1 (equal) & PASS \\
P1-COMP-01 & COMP-H2 (equal) & PASS \\
P1-COMP-01 & COMP-H3 (equal) & PASS \\
P1-COMP-01 & COMP-H4 (equal) & PASS \\
P1-COMP-01 & COMP-H5 (equal) & PASS \\
\bottomrule
\end{tabular}
\end{table}


\section{P1-REAL-01 details}
\label{app:real}

\emph{Pinned code facts (read, then executed through the adapter):} the step is \texttt{step\_scale * exp(log\_step)} with a per-mode timescale; the ZOH input matrix is \texttt{(1/Lambda)*(Lambda\_bar-1)*B} without \texttt{expm1}; the scan is $x_k=\bar\Lambda_kx_{k-1}+\bar B_ku_k$ with the gap before step $k$; each block applies \texttt{BatchNorm1d(affine=False)} over batch and time (running statistics at evaluation); an optional causal convolution defaults to off; the classifier head pools with \texttt{h.mean(dim=1)}.
\emph{Adapter check (untrained, four training trajectories):} forward outputs finite; gradients finite for all 54 parameters that receive one (the six static $\mathrm{Re}\,\Lambda$, $B$, $C$ tensors are unused in input-dependent mode; parameters upstream of the zero-initialized head projections receive zero gradient at step~0 only); an Adam step changed the parameters and lowered the loss on the same batch; evaluation was deterministic, and a per-step step tensor of ones matched the scalar step exactly.
\emph{Labels:} teacher output on the continuous path from an independent RK4 integration (128 sub-steps per training interval); the pooled label is its time average by Simpson's rule.
\emph{Readout of J1:} the last native output at or before each training-grid time.

\section{Numerical floor}
\label{app:numerics}

The naive coefficient $(e^{z}-1)/\lambda$ loses relative precision as $|z|$ shrinks ($\NoneNaive$ at $z{=}-10^{-14}$ in float64), while \texttt{expm1} stays at rounding level (\cref{tab:numerics}).
In emulated float32, split composition drifts to $\NfourFloor$ at $16{,}384$ sub-steps, so differences below about $10^{-4}$ are not interpretable at that length.

% AUTO-GENERATED by scripts/make_paper_assets.py from results/raw -- do not edit by hand.
\begin{table}[h]
\caption{Numerical floor (FR-E4). Left: relative error of the ZOH input coefficient $(e^{z}-1)/\lambda$ against a 50-digit reference. Right: splitting $[0,1]$ into $m$ held sub-steps
($\lambda{=}-1$, $g{=}\ln2$, exact $h(1)=0.5$); float32 is an \emph{emulation} (every primitive rounded, correctly rounded $\exp$), not a GPU kernel.}
\label{tab:numerics}
\centering\small
\begin{tabular}{rcc}
\toprule
$z$ & naive & \texttt{expm1} \\
\midrule
$-0.01$ & $5.4\times10^{-15}$ & $4.7\times10^{-17}$ \\
$-1\times10^{-4}$ & $1.3\times10^{-14}$ & $3.4\times10^{-17}$ \\
$-1\times10^{-6}$ & $1.6\times10^{-11}$ & $7.9\times10^{-17}$ \\
$-1\times10^{-8}$ & $1.1\times10^{-9}$ & $6.4\times10^{-17}$ \\
$-1\times10^{-10}$ & $8.3\times10^{-8}$ & $3.4\times10^{-17}$ \\
$-1\times10^{-12}$ & $2.2\times10^{-5}$ & $2.4\times10^{-17}$ \\
$-1\times10^{-14}$ & $8.0\times10^{-4}$ & $4.9\times10^{-17}$ \\
\bottomrule
\end{tabular}\hspace{1em}
\begin{tabular}{rccc}
\toprule
$m$ & fp64 \texttt{expm1} & fp32 \texttt{expm1} & fp32 naive \\
\midrule
$1$ & $0$ & $0$ & $0$ \\
$4$ & $2.2\times10^{-16}$ & $0$ & $6.0\times10^{-8}$ \\
$16$ & $2.2\times10^{-16}$ & $1.2\times10^{-7}$ & $2.4\times10^{-7}$ \\
$64$ & $4.4\times10^{-16}$ & $6.0\times10^{-8}$ & $0$ \\
$256$ & $1.8\times10^{-15}$ & $4.3\times10^{-6}$ & $5.0\times10^{-6}$ \\
$1024$ & $4.4\times10^{-15}$ & $1.8\times10^{-5}$ & $2.0\times10^{-5}$ \\
$4096$ & $1.1\times10^{-13}$ & $2.1\times10^{-5}$ & $2.8\times10^{-5}$ \\
$16384$ & $5.2\times10^{-13}$ & $1.1\times10^{-4}$ & $2.3\times10^{-4}$ \\
\bottomrule
\end{tabular}
\end{table}


\section{Additional toy tables}
\label{app:toy}

% AUTO-GENERATED by scripts/make_paper_assets.py from results/raw -- do not edit by hand.
\begin{table}[t]
\caption{Stiff scalar mode (FR-E4-N2; $g{=}\ln 2$, base step $0.5$, held sine input). Bilinear is stable ($|\bar A|<1$) for every step, but for $z{=}\Delta\lambda<-2$ its transition is negative
(sign-alternating transients) and tends to $-1$, whereas the exact decay $e^{z}$ tends to $0$. The last two columns are the output artifact against exact ZOH.
The stable/sign/decay columns are an exploratory re-reading of logged values.}
\label{tab:stiff}
\centering\small
\begin{tabular}{rrrrrccrr}
\toprule
$\lambda$ & $m$ & $z$ & $\bar A_{\mathrm{bil}}$ & $e^{z}$ & stable & sign-alt. & bilinear & Euler-$B$ \\
\midrule
$-5$ & $1$ & $-1.73$ & $+0.072$ & $0.18$ & yes & no & $0.098$ & $1.1$ \\
$-5$ & $4$ & $-0.43$ & $+0.644$ & $0.65$ & yes & no & $4.7\times10^{-3}$ & $0.23$ \\
$-5$ & $16$ & $-0.11$ & $+0.897$ & $0.90$ & yes & no & $2.9\times10^{-4}$ & $0.055$ \\
$-5$ & $32$ & $-0.05$ & $+0.947$ & $0.95$ & yes & no & $7.3\times10^{-5}$ & $0.027$ \\
$-50$ & $1$ & $-17.33$ & $-0.793$ & $3.0\times10^{-8}$ & yes & yes & $0.54$ & $16$ \\
$-50$ & $4$ & $-4.33$ & $-0.368$ & $0.013$ & yes & yes & $0.015$ & $3.4$ \\
$-50$ & $16$ & $-1.08$ & $+0.297$ & $0.34$ & yes & no & $2.1\times10^{-8}$ & $0.64$ \\
$-50$ & $32$ & $-0.54$ & $+0.574$ & $0.58$ & yes & no & $8.5\times10^{-9}$ & $0.30$ \\
$-500$ & $1$ & $-173.29$ & $-0.977$ & $5.5\times10^{-76}$ & yes & yes & $1.1$ & $170$ \\
$-500$ & $4$ & $-43.32$ & $-0.912$ & $1.5\times10^{-19}$ & yes & yes & $0.62$ & $42$ \\
$-500$ & $16$ & $-10.83$ & $-0.688$ & $2.0\times10^{-5}$ & yes & yes & $2.0\times10^{-3}$ & $9.8$ \\
$-500$ & $32$ & $-5.42$ & $-0.461$ & $4.4\times10^{-3}$ & yes & yes & $1.3\times10^{-11}$ & $4.4$ \\
\bottomrule
\end{tabular}
\vspace{4pt}

\begin{tabular}{rccc}
\toprule
$dt$ & exact ZOH & Euler-$B$ & bilinear \\
\midrule
$0.01$ & $8.0\times10^{-15}$ & $3.5\times10^{-3}$ & $0$ \\
$0.5$ & $0$ & $0.18$ & $1.1\times10^{-16}$ \\
$8$ & $0$ & $4.6$ & $5.6\times10^{-6}$ \\
\bottomrule
\end{tabular}
\vspace{2pt}

{\footnotesize Lower block (FR-E4-N3): relative steady-state error under a held unit input, $\lambda{=}-1$.}
\end{table}

% AUTO-GENERATED by scripts/make_paper_assets.py from results/raw -- do not edit by hand.
\begin{table}[t]
\caption{Adding new observations of the same signal (FR-E2-NEWOBS; toy, 32 units). The change $T=y_m-y_1$ at base times splits \emph{exactly} as $T=P+R$, where $P$ is the change of the continuous-time solution on the new held path
(the change in the held-input reconstruction; no task or Bayesian information gain is implied) and $R$ is the change of the discretization artifact. Hence $\|T\|^2=\|P\|^2+\|R\|^2+2\langle P,R\rangle$; we report these three terms as fractions of $\sum_{\text{units}}\|T\|^2$ (they sum to 1)
and the per-unit range of the cross fraction. $P$ uses the exact-ZOH output as a proxy for the continuous-time held-path solution (validated against RK4 to the bound shown).
Median total change is relative to $\|y_1\|$. EXPLORATORY re-aggregation; the cross term shows that no per-unit ``share'' is well defined, and none of these fractions is a causal contribution.}
\label{tab:newobs}
\centering\small
\begin{tabular}{lcccccc}
\toprule
variant & $m$ & median $\|T\|/\|y_1\|$ & $\|P\|^2$ & $\|R\|^2$ & $2\langle P,R\rangle$ & cross range \\
\midrule
exact ZOH, $\Delta{=}dt\,g$ & $2$ & $0.11$ & \multicolumn{4}{c}{artifact $\le1.7\times10^{-13}$ vs.\ RK4 (proxy reference)} \\
exact ZOH, $\Delta{=}dt\,g$ & $8$ & $0.19$ & \multicolumn{4}{c}{artifact $\le1.7\times10^{-13}$ vs.\ RK4 (proxy reference)} \\
Euler-$B$, $\Delta{=}dt\,g$ & $2$ & $0.15$ & $0.173$ & $0.567$ & $+0.260$ & $[-0.93,+0.42]$ \\
Euler-$B$, $\Delta{=}dt\,g$ & $8$ & $0.25$ & $0.198$ & $0.533$ & $+0.269$ & $[-1.37,+0.42]$ \\
bilinear, $\Delta{=}dt\,g$ & $2$ & $0.11$ & $0.849$ & $0.025$ & $+0.126$ & $[-1.04,+0.34]$ \\
bilinear, $\Delta{=}dt\,g$ & $8$ & $0.19$ & $0.886$ & $0.013$ & $+0.101$ & $[-0.69,+0.29]$ \\
exact ZOH, $\Delta{=}\tau g$ & $2$ & $0.42$ & $0.071$ & $1.260$ & $-0.331$ & $[-1.06,+0.34]$ \\
exact ZOH, $\Delta{=}\tau g$ & $8$ & $1.1$ & $0.025$ & $1.223$ & $-0.248$ & $[-1.02,+0.03]$ \\
Euler-$B$, $\Delta{=}\tau g$ & $2$ & $0.42$ & $0.053$ & $1.225$ & $-0.278$ & $[-0.72,+0.32]$ \\
Euler-$B$, $\Delta{=}\tau g$ & $8$ & $1.1$ & $0.020$ & $1.202$ & $-0.222$ & $[-0.75,+0.04]$ \\
bilinear, $\Delta{=}\tau g$ & $2$ & $0.42$ & $0.072$ & $1.260$ & $-0.331$ & $[-1.10,+0.35]$ \\
bilinear, $\Delta{=}\tau g$ & $8$ & $1.1$ & $0.025$ & $1.222$ & $-0.247$ & $[-1.04,+0.03]$ \\
\bottomrule
\end{tabular}
\end{table}

% AUTO-GENERATED by scripts/make_paper_assets.py from results/raw -- do not edit by hand.
\begin{table}[t]
\caption{Readouts under a density change (FR-E3-READOUT; toy, 32 units). Columns 2--3: scale-normalized change $|r_8-r_1|/s$ when only the first half of the base intervals is split into $8$ held sub-steps
(physical path unchanged); $s$ is the RMS base-time output at $m{=}1$ ($\times K_1$ for the sample sum). Column 4: scale-normalized error to the continuous-time truth when new observations are added to the first half only, $m{=}1\to8$.
Medians over units. The exact held-path integral is defined only for exact-ZOH states.}
\label{tab:readout}
\centering\small
\begin{tabular}{lccc}
\toprule
readout & split, exact ZOH $\Delta{=}dt\,g$ & split, Euler-$B$ $\Delta{=}\tau g$ & new obs., exact ZOH \\
\midrule
last sample & $8.9\times10^{-17}$ & $0.19$ & $0.085$$\to$$0.099$ \\
sample sum & $1.5$ & $2.9$ & -- \\
sample mean & $0.30$ & $0.46$ & $0.060$$\to$$0.28$ \\
right-endpoint $\sum y_k\,dt_k$ & $0.026$ & $0.24$ & $0.060$$\to$$0.043$ \\
trapezoidal endpoint sum & $0.011$ & $0.25$ & $0.039$$\to$$0.023$ \\
exact held-path integral & $1.5\times10^{-16}$ & n/a & $0.039$$\to$$0.032$ \\
\bottomrule
\end{tabular}
\end{table}


\section{Proof sketches (unproved)}
\label{app:proofs}

\paragraph{Assumptions.}
(A1) $\mathrm{Re}\,\lambda_n\le-\mu<0$; (A2) $g:[-U,U]\to[g_{\min},g_{\max}]$, $g_{\min}>0$, Lipschitz $L_g$; (A3) $F(v)=g(v)B(v)v$ is $L_F$-Lipschitz with $|F|\le F_{\max}$; (A4) $u$ is $L_u$-Lipschitz with $|u|\le U$.

\paragraph{Hold error (UNPROVED sketch).}
For solutions driven by the held and the continuous input, $\dot e=g(\tilde u)\lambda e+[(g(\tilde u)-g(u))\lambda h_u+F(\tilde u)-F(u)]$; the bracket is at most $(L_g|\lambda|H+L_F)L_u\delta$ with $H=F_{\max}/(g_{\min}\mu)$ and the propagator is bounded by $e^{-g_{\min}\mu(t-s)}$, suggesting $\sup_t|e(t)|\le(L_g|\lambda|H+L_F)L_u\delta/(g_{\min}\mu)$; a complete proof must treat the discontinuous $\tilde u$ and has not been written.

\paragraph{Euler-$B$ (UNPROVED sketch).}
$e^z-1-z=z^2\int_0^1(1-s)e^{sz}ds$ gives the local bound; summing against the contraction suggests a global $O(\delta)$ bound whose constant grows with $\Delta|\lambda|$.

\paragraph{Refinement with $\Delta{=}\tau g$ (CONJECTURE).}
The recurrence is the exact flow of $\dot h=(\tau/\delta)g(\tilde u)(\lambda h+B(\tilde u)\tilde u)$, whose solution should approach $-B(u)u/\lambda$ as $\delta\to0$, so the artifact would tend to a nonzero limit.

\section{Reproducibility}
\label{app:repro}

The anonymized code contains the toy, the RK4 references, the cascade check, the TIDES adapter, convention check, runner and diagnostic, the frozen checkpoint with its hash, and the generator that produces every table, figure and number in this paper from the raw records.

\end{document}
